\begin{initiativkarte}{SAIL}

  % Bild (optional)
  \IfFileExists{bilder/sail.jpg}{%
    \begin{center}
      \includegraphics[width=\linewidth,height=5cm,keepaspectratio]{bilder/sail.jpg}
    \end{center}
    \vspace{0.4cm}
  }{}

  % Kurzbeschreibung
  \textit{\textcolor{hawgray}{SAIL ist das Open-Access-Publikationsportal der HAW Hamburg für studentische Arbeiten und ermöglicht freien Zugang zu wissenschaftlichem Wissen.}}

  \vspace{0.5cm}
  \hawrule

  % Beschreibung
  \textbf{\textcolor{hawblue}{Was ist SAIL?}}\\[0.3cm]

  SAIL ist das Open-Access-Publikationsportal der Hochschule für Angewandte Wissenschaften Hamburg.

  Studierende können dort ihre im Studium entstandenen Arbeiten veröffentlichen und frei zugänglich machen.

  Ziel des Portals ist es, Wissen offen zu teilen und eine freie, uneingeschränkte Nutzung von studentischen Arbeiten zu ermöglichen.

  \vspace{0.3cm}

  Alle Publikationen werden inhaltlich erschlossen, langfristig archiviert und durch Lehrende oder Peer-Review-Verfahren geprüft, um eine hohe Qualität sicherzustellen.

  \vspace{0.5cm}

  \textbf{Funktionen von SAIL}
  \begin{itemize}
    \item Veröffentlichung studentischer Arbeiten
    \item Freier Zugang zu wissenschaftlichen Texten
    \item Recherche in bestehenden Publikationen
    \item Langzeitarchivierung von Arbeiten
    \item Qualitätssicherung durch Begutachtung
  \end{itemize}

  \vspace{0.5cm}

  Das Portal befindet sich derzeit noch im Aufbau und wurde im Wintersemester 2020/21 von Studierenden des Departments Informatik entwickelt.

  Betrieben wird SAIL vom Open-Access-Lab der HAW Hamburg unter Leitung von Prof. Dr. Ulrike Verch und Dr. Steffen Rudolph.

  \vspace{0.5cm}

  % Tags
  \tag{Open Access}
  \tag{Wissenschaft}
  \tag{Publikation}
  \tag{Forschung}
  \tag{Wissen teilen}

  \vspace{0.6cm}
  \hawrule

  % Infozeilen
  \infozeile{\faUsers}{Zielgruppe: Alle Studierenden der HAW Hamburg}
  \infozeile{\faGlobe}{Website: Open-Access-Lab HAW Hamburg}
  \infozeile{\faBook}{Fokus: Veröffentlichung und freie Nutzung studentischer Arbeiten}

\end{initiativkarte}